\section{Opis algorytmów}
\label{sec:algorytmy}

W~ramach projektu rozważano i~częściowo zaimplementowano kilka podejść do
generowania trasy. Głównymi meta-heurystykami, które zostały zaimplementowane
i~służą do generowania populacji rozwiązań dopuszczalnych, są \emph{Artificial
Bee Colony} (ABC) oraz \emph{Bee Algorithm} (BA, tzw.\ taniec pszczół). Dodatkowo,
jako prosty baseline porównawczy i~mechanizm testowania interfejsu,
zaimplementowano dwie pomocnicze procedury:
\emph{Random Walk} oraz \emph{Greedy attractions} (zachłanne dążenie do najbliższej
atrakcji).

\subsection{Artificial Bee Colony (ABC)}
\label{sec:abc}

ABC to populacyjna meta-heurystyka inspirowana zachowaniem pszczół. W~klasycznej
formie wyróżnia trzy typy ,,pszczół'':
\begin{itemize}
    \item \textbf{Pszczoły zatrudnione} (\emph{employed bees}) --- przypisane
          do konkretnych rozwiązań (źródeł pokarmu) i~modyfikujące je lokalnie.
    \item \textbf{Obserwatorzy} (\emph{onlookers}) --- wybierają źródło
          z~prawdopodobieństwem rosnącym z~jakością źródła i~również je modyfikują.
    \item \textbf{Zwiadowcy} (\emph{scouts}) --- zastępują źródła, które od~zbyt
          wielu prób (\emph{trial}) nie poprawiły się, nowymi losowymi rozwiązaniami.
\end{itemize}

Schemat iteracji ABC dla naszego problemu przedstawia~\Cref{alg:abc-overview}.

\begin{algorithm}[H]
\caption{Generowanie populacji metodą ABC (ogólny zarys)}
\label{alg:abc-overview}
\begin{algorithmic}[1]
\Require dane wejściowe \texttt{parsed\_data}, rozmiar populacji $N$, liczba iteracji $I$, limit stagnacji $L$, liczba onlookerów $O$
\Ensure populacja $\mathcal{P} = \{s_1, \dots, s_N\}$ rozwiązań dopuszczalnych
\State $\mathcal{P} \gets \emptyset$
\While{$|\mathcal{P}| < N$}
    \State $s \gets$ \Call{LosoweRozwiazanieDopuszczalne}{}
    \If{$s \ne \bot$ \textbf{i} $s$ nie jest duplikatem w~$\mathcal{P}$}
        \State $\mathcal{P} \gets \mathcal{P} \cup \{s\}$
    \EndIf
\EndWhile
\For{$it = 1, \dots, I$}
    \For{$i = 1, \dots, N$} \Comment{Faza employed}
        \State $s' \gets$ \Call{MutacjaWaypointow}{$\mathcal{P}_i$}
        \If{$s' \ne \bot$ \textbf{i} $s' \succ \mathcal{P}_i$}
            \State $\mathcal{P}_i \gets s'$;\ \  $\mathcal{P}_i.\text{trial} \gets 0$
        \Else
            \State $\mathcal{P}_i.\text{trial} \gets \mathcal{P}_i.\text{trial} + 1$
        \EndIf
    \EndFor
    \State $w \gets$ \Call{WagiRankingowe}{$\mathcal{P}$}
    \For{$k = 1, \dots, O$} \Comment{Faza onlooker}
        \State $i \gets$ \Call{WylosujZWagami}{$w$}
        \State $s' \gets$ \Call{MutacjaWaypointow}{$\mathcal{P}_i$}
        \If{$s' \ne \bot$ \textbf{i} $s' \succ \mathcal{P}_i$}
            \State $\mathcal{P}_i \gets s'$;\ \  $\mathcal{P}_i.\text{trial} \gets 0$
        \Else
            \State $\mathcal{P}_i.\text{trial} \gets \mathcal{P}_i.\text{trial} + 1$
        \EndIf
    \EndFor
    \For{$i = 1, \dots, N$} \Comment{Faza scout}
        \If{$\mathcal{P}_i.\text{trial} \ge L$}
            \State $\mathcal{P}_i \gets$ \Call{LosoweRozwiazanieDopuszczalne}{}
            \State $\mathcal{P}_i.\text{trial} \gets 0$
        \EndIf
    \EndFor
\EndFor
\State \Return $\mathcal{P}$
\end{algorithmic}
\end{algorithm}

Relacja porządku $s' \succ s$ to leksykograficzne porównanie krotek:
\[
\big(\,\mathrm{Value}(s'),\ -\mathrm{Time}(s'),\ -\mathrm{Cost}_{\mathrm{ruch}}(s'),\ -\mathrm{Cost}_{\mathrm{atrakcji}}(s')\,\big)
\;>\;
\big(\,\mathrm{Value}(s),\ -\mathrm{Time}(s),\ -\mathrm{Cost}_{\mathrm{ruch}}(s),\ -\mathrm{Cost}_{\mathrm{atrakcji}}(s)\,\big).
\]

\subsubsection{Reprezentacja rozwiązania}
\label{sec:repr}

Rozwiązanie w~populacji jest dwupoziomowe:

\begin{enumerate}
    \item \textbf{Plan logiczny} --- lista \emph{waypointów}
          $\mathbf{W} = \langle s, a_{i_1}, a_{i_2}, \dots, a_{i_k}, e \rangle$,
          gdzie $s,e$ to punkty startu i~końca, a~$a_{i_j}$ to wybrane atrakcje
          w~zaplanowanej kolejności odwiedzenia.
    \item \textbf{Realizacja fizyczna} --- konkretna ścieżka komórek
          $P = \langle p_0, p_1, \dots, p_L \rangle$, zbudowana z~planu
          logicznego przez łączenie kolejnych waypointów najtańszą ścieżką
          A\textsuperscript{*} z~uwzględnieniem zakazu używania już wykorzystanych
          krawędzi i~niezaplanowanych komórek atrakcji.
\end{enumerate}

Każde rozwiązanie w~populacji przechowuje również:
\begin{itemize}
    \item \texttt{path} --- pełną listę komórek (lista par $(r,c)$),
    \item \texttt{visited\_attractions} --- listę odwiedzonych atrakcji,
    \item \texttt{movement\_cost}, \texttt{attraction\_cost}, \texttt{total\_time}, \texttt{total\_value},
    \item \texttt{type\_counts} --- wektor liczności atrakcji każdego typu,
    \item flagi dopuszczalności: \texttt{budget\_ok}, \texttt{time\_ok}, \texttt{type\_ok}, \texttt{feasible},
    \item \texttt{trial} --- licznik nieudanych modyfikacji (na potrzeby zwiadowców),
    \item \texttt{id}, \texttt{waypoints} --- identyfikator i~plan logiczny.
\end{itemize}

\subsubsection{Wygenerowanie populacji początkowej}
\label{sec:init-pop}

Populacja początkowa składa się z~$N$ losowych rozwiązań \textbf{dopuszczalnych}.
Procedura konstrukcji pojedynczego rozwiązania składa się z~dwóch kroków:

\begin{enumerate}
    \item Losowy wybór listy waypointów spełniający ograniczenia liczności
          typów (\Cref{alg:random-waypoints}).
    \item Materializacja ścieżki przez łączenie kolejnych waypointów
          algorytmem A\textsuperscript{*} z~zakazem powtarzania krawędzi
          i~komórek (\Cref{alg:build-path}). Jeżeli ścieżka się nie udaje
          (np.\ nie ma fizycznej drogi spełniającej ograniczenia),
          rozwiązanie jest odrzucane i~próbujemy następne.
\end{enumerate}

\begin{algorithm}[H]
\caption{\textsc{LosoweWaypointy} --- losowanie planu atrakcji}
\label{alg:random-waypoints}
\begin{algorithmic}[1]
\Require dane mapy, punkty $s, e$, parametr $E_{\max}$ (maksymalna liczba dodatkowych atrakcji)
\Ensure lista waypointów $\mathbf{W}$ lub $\bot$ przy braku spełnienia minimów
\State $\mathrm{chosen} \gets \emptyset$
\ForAll{typ $t$}
    \State $n_t \gets m_t -$ liczba atrakcji typu $t$ obecnych w~$\{s, e\}$
    \If{$n_t > 0$}
        \State wylosuj $n_t$ atrakcji typu $t$ spośród dostępnych
        \If{kandydatów jest za mało} \State \Return $\bot$ \EndIf
        \State dołącz wylosowane do $\mathrm{chosen}$
    \EndIf
\EndFor
\State $E \gets \mathrm{rand}\!\left(0,\ \min(E_{\max},\ |\text{wolnych atrakcji}|)\right)$
\For{$k = 1, \dots, E$}
    \State wylosuj atrakcję $a$ z~prawdopodobieństwem ważonym
           $\frac{v_a + 1}{\tau_{t_a} + 1}$
    \If{$a \notin \mathrm{chosen}$ \textbf{i} dorzucenie nie łamie $M_{t_a}$}
        \State $\mathrm{chosen} \gets \mathrm{chosen} \cup \{a\}$
    \EndIf
\EndFor
\State pomieszaj kolejność $\mathrm{chosen}$
\State \Return konkatenacja $\langle s \rangle$, $\mathrm{chosen}$, $\langle e \rangle$
\end{algorithmic}
\end{algorithm}

\begin{algorithm}[H]
\caption{\textsc{BudujSciezkeZWaypointow} --- konstrukcja ścieżki fizycznej}
\label{alg:build-path}
\begin{algorithmic}[1]
\Require mapa, waypointy $\mathbf{W} = \langle w_0, w_1, \dots, w_K \rangle$
\Ensure ścieżka $P$ (lista komórek) lub $\bot$ jeśli nie istnieje
\State $\mathrm{visited} \gets \{w_0\};\quad \mathrm{usedEdges} \gets \emptyset$
\State $P \gets \langle w_0 \rangle$
\For{$i = 0, \dots, K-1$}
    \State $\text{blokady} \gets (\mathcal{A}\text{-pozycje} \setminus \{w_0, w_K, w_i, w_{i+1}\})\ \cup\ (\mathrm{visited} \setminus \{w_i, w_{i+1}\})$
    \State $S \gets$ \Call{AStar}{$w_i, w_{i+1}$, blokady, $\mathrm{usedEdges}$}
    \If{$S = \bot$} \Return $\bot$ \EndIf
    \For{każda kolejna krawędź $(u,v)$ w~$S$}
        \State $\mathrm{usedEdges} \gets \mathrm{usedEdges} \cup \{\{u,v\}\}$
    \EndFor
    \For{każdy kolejny węzeł $p$ w~$S$ poza pierwszym}
        \If{$p \in \mathrm{visited}$} \Return $\bot$ \EndIf
        \State $\mathrm{visited} \gets \mathrm{visited} \cup \{p\}$;\ \  dołącz $p$ do~$P$
    \EndFor
\EndFor
\State \Return $P$
\end{algorithmic}
\end{algorithm}

Procedura A\textsuperscript{*} jest standardowa --- używa heurystyki
optymistycznej $h(p, g) = w_{\min} \cdot (1{,}4\cdot d_{\mathrm{diag}} +
d_{\mathrm{ax}})$, gdzie $w_{\min}$ to minimalna waga na siatce,
$d_{\mathrm{diag}}$ to liczba kroków diagonalnych, a~$d_{\mathrm{ax}}$ to
liczba kroków osiowych w~ruchu Chebysheva. Zapewnia to dopuszczalność
heurystyki i~optymalność znalezionej trasy w~sensie kosztu ruchu (przy
zadanych blokadach komórek i~zakazach krawędzi).

\subsubsection{Operatory modyfikacji rozwiązania (sąsiedztwo)}
\label{sec:operators}

Modyfikacja rozwiązania (\textsc{MutacjaWaypointow}) polega na zmianie
listy waypointów, a~następnie próbie materializacji nowej ścieżki.
Wybór operatora jest losowy z~rozkładem przedstawionym
w~\Cref{tab:operators}. W~każdej próbie sprawdzane są ograniczenia
liczności typów; jeżeli kandydat nie da się zmaterializować lub nie spełnia
ograniczeń, próba kończy się fiaskiem.

\begin{table}[H]
\centering
\caption{Operatory mutacji listy waypointów (sąsiedztwo).}
\label{tab:operators}
\begin{tabular}{p{3.3cm} c p{8.3cm}}
\toprule
\textbf{Operator} & \textbf{P} & \textbf{Opis} \\
\midrule
Dodaj atrakcję      & 0{,}30 & Wybierany kandydat $a \notin \mathbf{W}$
                                ze zbioru 20 najlepiej rokujących wg
                                $(v_a + 1) / (\tau_{t_a} + 1)$; wstawiany
                                w~losowej pozycji, jeśli nie łamie $M_{t_a}$. \\
Usuń atrakcję       & 0{,}25 & Losowo wybierany element $\mathbf{W}$, którego
                                usunięcie nie złamie $m_{t}$. \\
Zamień sąsiadujące  & 0{,}20 & Wybierane dwie pozycje w~$\mathbf{W}$
                                i~zamieniane miejscami. \\
2-opt (odwróć fragment) & 0{,}20 & Losowo wybierany fragment $[i, j]$
                                i~odwracany. \\
Przesuń atrakcję    & 0{,}05 & Wyjęcie pojedynczego elementu i~wstawienie
                                w~innym miejscu (\emph{shift/relocate}). \\
\bottomrule
\end{tabular}
\end{table}

\subsubsection{Selekcja typu onlooker (wagi rankingowe)}
\label{sec:onlookers}

W~fazie onlookerów wybór modyfikowanego źródła odbywa się na podstawie
\emph{rangi} rozwiązania w~populacji (a~nie surowych wartości funkcji celu),
co zapobiega dominacji jednego źródła. Procedura wag rankingowych jest
następująca:

\begin{enumerate}
    \item Posortuj populację malejąco wg krotki porządkowej (patrz
          \Cref{sec:abc}).
    \item Rozwiązaniu o~randze $r$ (od~0) przypisz wagę $w_r = N - r$.
    \item Losuj indeks źródła z~rozkładem proporcjonalnym do $w_r$.
\end{enumerate}

Liczba prób onlookerów w~iteracji wynosi $O$ (parametr; domyślnie
$O = N$).

\subsubsection{Faza zwiadowcza (scout)}
\label{sec:scouts}

Jeśli licznik nieudanych prób ($\mathrm{trial}$) dla rozwiązania przekroczy
limit $L$, rozwiązanie jest zastępowane nowym losowym rozwiązaniem
dopuszczalnym (\Cref{alg:random-waypoints} + \Cref{alg:build-path}). Po
udanym zastąpieniu licznik prób się zeruje. Jeżeli nie udaje się
wygenerować nowego rozwiązania w~zadanej liczbie prób, licznik również się
zeruje, by nie zapętlać próby zastąpienia.

\subsubsection{Pseudokod operatora mutacji}
\label{sec:mutate-pseudo}

\begin{algorithm}[H]
\caption{\textsc{MutacjaWaypointow}}
\label{alg:mutate}
\begin{algorithmic}[1]
\Require rozwiązanie $s$, $\mathrm{maxAttempts}$
\Ensure rozwiązanie $s'$ (jeśli udało się znaleźć dopuszczalnego sąsiada) lub $\bot$
\For{$j = 1, \dots, \mathrm{maxAttempts}$}
    \State $\mathbf{W}' \gets s.\mathbf{W}$
    \State $\rho \gets \mathrm{rand}(0, 1)$
    \If{$\rho < 0{,}30$} \Comment{Dodaj}
        \State dobierz kandydata $a$ ze zbioru top-20 wg $(v_a+1)/(\tau_{t_a}+1)$
        \If{liczność typu $t_a$ po dodaniu $\le M_{t_a}$}
            \State wstaw $a$ w~losowe miejsce $\mathbf{W}'$
        \Else \State \textbf{continue}
        \EndIf
    \ElsIf{$\rho < 0{,}55$} \Comment{Usuń}
        \State $R \gets$ pozycje, których usunięcie nie złamie $m_{t}$
        \If{$R = \emptyset$} \textbf{continue} \EndIf
        \State usuń element o~losowej pozycji z~$R$
    \ElsIf{$\rho < 0{,}75$} \Comment{Zamień}
        \If{$|\mathbf{W}'| < 2$} \textbf{continue} \EndIf
        \State wylosuj $i, k$; zamień $\mathbf{W}'[i] \leftrightarrow \mathbf{W}'[k]$
    \ElsIf{$\rho < 0{,}95$} \Comment{2-opt}
        \If{$|\mathbf{W}'| < 4$} \textbf{continue} \EndIf
        \State wylosuj $i \le k$; odwróć $\mathbf{W}'[i..k]$
    \Else \Comment{Przesuń}
        \If{$|\mathbf{W}'| < 2$} \textbf{continue} \EndIf
        \State wylosuj $i$, $k$; przenieś $\mathbf{W}'[i]$ na pozycję $k$
    \EndIf
    \State $P \gets$ \Call{BudujSciezkeZWaypointow}{$\mathbf{W}'$}
    \If{$P \ne \bot$ \textbf{i} ocena $P$ daje rozwiązanie dopuszczalne}
        \State \Return rozwiązanie zbudowane z~$\mathbf{W}'$ i~$P$
    \EndIf
\EndFor
\State \Return $\bot$
\end{algorithmic}
\end{algorithm}

\subsubsection{Funkcja oceny ścieżki}
\label{sec:eval}

Funkcja \texttt{evaluate\_path} (plik \texttt{initial\_population.py}) liczy:
\begin{itemize}
    \item \texttt{movement\_cost} --- $\sum_{l} w_{p_l}\cdot\delta_l$,
    \item \texttt{attraction\_cost} --- suma $k_i$ po unikalnych odwiedzonych atrakcjach,
    \item \texttt{total\_time} --- $\texttt{movement\_cost} + \sum \tau_{t_i}$,
    \item \texttt{total\_value} --- $\sum v_i$ po unikalnych odwiedzonych atrakcjach,
    \item \texttt{type\_counts} --- wektor $\big[\sum_{i: t_i = t} \mathbb{1}_{a_i \in \mathcal{V}(P)}\big]_{t = 0, \dots, T-1}$,
    \item flagi: \texttt{budget\_ok} ($\le B$), \texttt{time\_ok} ($\le T_{\max}$),
          \texttt{type\_ok} ($m_t \le \cdot \le M_t$) i~\texttt{feasible} (koniunkcja).
\end{itemize}

Spójność z~modelem matematycznym jest zapewniona: zwracana krotka cech
odpowiada wprost ograniczeniom \eqref{eq:budget}--\eqref{eq:types}.

\subsection{Bee Algorithm (BA, taniec pszczół)}
\label{sec:bee}

Bee Algorithm (BA) to populacyjna meta-heurystyka inspirowana sposobem, w~jaki
pszczoły \emph{tańcem} komunikują jakość znalezionych źródeł pokarmu i~rekrutują
inne osobniki do dalszej eksploracji otoczenia. W~naszej implementacji BA działa
na \textbf{tej samej reprezentacji rozwiązania} (waypointy + materializacja
A\textsuperscript{*}) i~na \textbf{tym samym sąsiedztwie} (operatory mutacji) co
ABC, a~różni się przede wszystkim mechanizmem \emph{rekrutacji} oraz
\emph{porzucania} rozwiązań.

Intuicyjnie odpowiada to następującym rolom:
\begin{itemize}
    \item \textbf{Rekruterzy (tańczące pszczoły)} --- podzbiór najlepszych rozwiązań,
          których jakość jest komunikowana i~które przyciągają kolejne próby
          ulepszania.
    \item \textbf{Zbieraczki (foragers)} --- wykonują lokalne przeszukiwanie,
          generując sąsiadów (mutacje waypointów) wokół rozwiązań-rekruterów.
    \item \textbf{Zwiadowcy (scouts)} --- wprowadzają różnorodność, zastępując część
          słabych / zastanych rozwiązań nowymi losowymi rozwiązaniami dopuszczalnymi.
\end{itemize}

Schemat iteracji BA dla naszego problemu przedstawia~\Cref{alg:bee-overview}.

\begin{algorithm}[H]
\caption{Generowanie populacji metodą BA (ogólny zarys)}
\label{alg:bee-overview}
\begin{algorithmic}[1]
\Require dane wejściowe \texttt{parsed\_data}, rozmiar populacji $N$, liczba iteracji $I$, prawdopodobieństwo rekrutacji $p_r$, udział zwiadowców $r_s$, próg jakości tańca $Q$ (\%), cierpliwość porzucenia $A$
\Ensure populacja $\mathcal{P} = \{s_1, \dots, s_N\}$ rozwiązań dopuszczalnych
\State $\mathcal{P} \gets$ $N$ losowych rozwiązań dopuszczalnych (por.\ \Cref{sec:init-pop})
\State $s^\star \gets \bot$
\For{$it = 1, \dots, I$}
    \State $b \gets \max\limits_{s \in \mathcal{P}} s$ \Comment{najlepsze wg relacji $\succ$}
    \If{$s^\star = \bot$ \textbf{lub} $b \succ s^\star$} \State $s^\star \gets b$ \EndIf
    \State $\theta \gets \mathrm{Value}(s^\star) \cdot (Q/100)$ \Comment{próg jakości tańca}
    \State $E \gets \{ i : \mathrm{Value}(\mathcal{P}_i) \ge \theta \}$ \Comment{rekruterzy}
    \If{$E = \emptyset$} \State $E \gets \{\arg\max_i \mathcal{P}_i\}$ \EndIf
    \State $w_i \gets \mathrm{Value}(\mathcal{P}_i) + 1$ dla $i \in E$ \Comment{wagi rekrutacji}
    \State $F \gets \max(1, \lfloor N \cdot p_r \rfloor)$ \Comment{liczba zbieraczek}
    \For{$k = 1, \dots, F$}
        \State $i \gets$ \Call{WylosujZWagami}{$E, w$}
        \State $s' \gets$ \Call{MutacjaWaypointow}{$\mathcal{P}_i$} \Comment{por.\ \Cref{sec:operators}}
        \If{$s' \ne \bot$ \textbf{i} $s' \succ \mathcal{P}_i$}
            \State $\mathcal{P}_i \gets s'$;\ \ $\mathcal{P}_i.\text{trial} \gets 0$
        \Else
            \State $\mathcal{P}_i.\text{trial} \gets \mathcal{P}_i.\text{trial} + 1$
        \EndIf
    \EndFor
    \State $S \gets \max(1, \lfloor N \cdot r_s \rfloor)$ \Comment{liczba zwiadowców}
    \State wybierz $S$ indeksów o~największym \texttt{trial}
    \ForAll{wybrane $i$}
        \State $\mathcal{P}_i \gets$ \Call{LosoweRozwiazanieDopuszczalne}{}
        \State $\mathcal{P}_i.\text{trial} \gets 0$
    \EndFor
    \For{$i = 1, \dots, N$} \Comment{porzucanie stagnujących źródeł}
        \If{$\mathcal{P}_i.\text{trial} \ge A$}
            \State $\mathcal{P}_i \gets$ \Call{LosoweRozwiazanieDopuszczalne}{}
            \State $\mathcal{P}_i.\text{trial} \gets 0$
        \EndIf
    \EndFor
\EndFor
\State \Return $\mathcal{P}$
\end{algorithmic}
\end{algorithm}

Relacja porządku $s' \succ s$ jest taka sama jak w~ABC (leksykograficzne
porównanie krotek jakości; por.\ \Cref{sec:abc}).

\subsubsection{Reprezentacja rozwiązania}

Reprezentacja (waypointy + materializacja ścieżki) jest identyczna jak w~ABC
(\Cref{sec:repr}).

\subsubsection{Wygenerowanie populacji początkowej}

Inicjalizacja populacji polega na~wygenerowaniu $N$ losowych rozwiązań
\textbf{dopuszczalnych} przez losowanie waypointów i~materializację ścieżki
(\Cref{sec:init-pop}, \Cref{alg:random-waypoints}, \Cref{alg:build-path}).

\subsubsection{Operatory modyfikacji (sąsiedztwo)}

Lokalne przeszukiwanie wykonywane przez zbieraczki wykorzystuje ten sam operator
\textsc{MutacjaWaypointow} co ABC (\Cref{sec:operators}, \Cref{tab:operators},
\Cref{alg:mutate}).

\subsubsection{Rekrutacja (taniec) i dobór obszarów eksploatacji}

W~każdej iteracji wyznaczany jest próg jakości tańca $Q$ (w~\%), a~następnie do
zbioru rekruterów trafiają te rozwiązania, których wartość
$\mathrm{Value}(s)$ jest nie mniejsza niż $Q\%$ najlepszej wartości
zaobserwowanej dotychczas. Zbieraczki losują następnie źródła do eksploracji
z~rozkładu proporcjonalnego do $\mathrm{Value}(s)+1$ (żeby uniknąć wag zerowych).
Parametr $p_r$ steruje liczbą prób ulepszania w~iteracji.

\subsubsection{Zwiadowcy i porzucanie rozwiązań}

W~implementacji zastosowano dwa mechanizmy dywersyfikacji:
\begin{enumerate}
    \item \textbf{Zwiadowcy} --- w~każdej iteracji część populacji (udział $r_s$)
          jest zastępowana nowymi losowymi rozwiązaniami dopuszczalnymi; jako
          kandydatów wybieramy te osobniki, które mają największy licznik
          stagnacji \texttt{trial}.
    \item \textbf{Porzucanie} --- jeżeli licznik \texttt{trial} przekroczy
          \texttt{abandon\_patience} $=A$, źródło jest uznawane za nieperspektywiczne
          i~również zastępowane nowym losowym rozwiązaniem.
\end{enumerate}

\subsection{Podejścia pomocnicze}
\label{sec:baselines}

W~kodzie znajdują się także uproszczone metody pomocnicze, używane głównie
do testowania interfejsu wizualizacji w~aplikacji Streamlit.

\subsubsection{Random Walk}

Algorytm losowy spaceruje od~$s$ po sąsiedztwie 8-kierunkowym, dopóki nie
zostanie wykorzystany budżet kosztu lub nie zostanie osiągnięty punkt $e$.
Nie pilnuje ograniczeń typów, nie unika powtarzania komórek/krawędzi
i~służy wyłącznie jako baseline. Implementacja: \texttt{src/path\_solver.py},
funkcja \texttt{solve\_random\_walk}.

\subsubsection{Greedy attractions}

W~każdym kroku wybierana jest najbliższa (w~metryce Chebysheva) jeszcze
nieodwiedzona atrakcja i~algorytm wykonuje kroki w~jej kierunku, dopóki
budżet pozwala. Implementacja: \texttt{src/path\_solver.py}, funkcja
\texttt{solve\_greedy\_attractions}. Także baseline.

\subsection{Inne pomysły rozważane w~fazie projektowej}
\label{sec:other-ideas}

W~trakcie projektu rozważano kilka alternatywnych dróg, które nie weszły
do finalnej implementacji:
\begin{itemize}
    \item \textbf{Algorytm genetyczny (GA)} --- z~populacją chromosomów
          jako permutacji wybranych atrakcji; krzyżowania typu PMX/OX,
          mutacje swap/insert/2-opt. Architektura aplikacji (\texttt{app.py})
          zakłada możliwość podłączenia GA przez stabilny interfejs
          \texttt{solve(map, params) → path}, więc rozszerzenie jest naturalne.
    \item \textbf{Programowanie liniowe całkowitoliczbowe (ILP)} --- model
          \eqref{eq:obj}--\eqref{eq:bin} można sprowadzić do mieszanej
          formulacji ILP z~przepływami w~grafie siatki; rozważano użycie
          solvera (np.\ CBC) do małych instancji jako wartość referencyjna.
          Decyzja: pominięto ze względu na koszt implementacji i~ograniczoną
          skalowalność.
    \item \textbf{Symulowane wyżarzanie (SA)} --- jako pojedyncze-rozwiązaniowa
          alternatywa, działająca na tej samej reprezentacji waypointów
          i~tym samym zbiorze operatorów. Pomysł porzucono, ponieważ ABC
          dawało już zadowalające wyniki, a~równoległa populacja jest
          wartością samą w~sobie (np.\ do dalszego ranking-u/eksploracji).
    \item \textbf{Ant Colony Optimization (ACO)} --- pomysł na klasycznej
          stronie konstrukcji ścieżki, jednak nieoczywista adaptacja do
          twardych ograniczeń typów; pominięto.
\end{itemize}

W~niniejszej dokumentacji szczegółowo opisane są algorytmy ABC oraz Bee Algorithm,
ponieważ oba podejścia zostały zaimplementowane i~wykorzystywane w~aplikacji.
